\thispagestyle{empty}
\vspace*{0.2cm}

\begin{center}
    \textbf{Zusammenfassung}
\end{center}

\vspace*{0.2cm}

\noindent 

\noindent
Die großräumige Überwachung der pflanzlichen Photosynthese ist wichtig, um die regionale Kohlenstoffaufnahme und landwirtschaftliche Produktivität zu bestimmen.
Die solarinduzierte Chlorophyllfluoreszenz (SIF) ist eine schwache Lichtemission des Chlorophylls während der Photosynthese und ein fernerkundlich beobachtbares Signal mit Bezug zur Bruttoprimärproduktion.
Der OCO-2-Satellit der NASA liefert derzeit die räumlich höchstaufgelösten SIF-Daten, jedoch mit lückenhafter räumlicher und zeitlicher Abdeckung.
Bestehende verbesserte SIF-Produkte erreichen lediglich etwa 500~m Auflösung und eignen sich daher nur eingeschränkt zur Isolierung kulturartenreiner Signale in kleinen, unregelmäßig geformten deutschen Ackerflächen.
Benötigt wird deshalb ein deutlich höher aufgelöstes SIF-Produkt auf Grundlage räumlich und zeitlich vollständiger Prädiktoren, insbesondere aus Sentinel-2-Daten.
Unklar ist, wie Modellarchitektur und Prädiktordarstellung als mittelwertaggregierte Tabellenwerte oder räumliche Bilddaten die Genauigkeit und räumliche Konsistenz innerhalb desselben Untersuchungsrahmens beeinflussen.
Zu untersuchen sind außerdem die Auswirkungen von Aggregation, Stichprobendichte, zeitlicher Abweichung zwischen SIF und Prädiktoren sowie agroklimatischer Variation.
Die Studie prüft zudem, ob kulturartenreines verbessertes SIF für die Winterweizen-Ertragsvorhersage nützlicher ist als SIF aus gemischten Footprints. Qualitätsgefilterte OCO-2-Messungen bei 757 und 771~nm wurden kombiniert und mithilfe von Sentinel-2-Spektralindizes mit 20~m Auflösung sowie Strahlungs-, Saison- und Kulturzusammensetzungsvariablen für Deutschland von Februar bis Juli 2019--2024 modelliert.
Verglichen wurden Random Forest, XGBoost und bildbasierte U-Nets.
Das beste U-Net wurde mit festen 4-km-Rasterzellen trainiert, die mindestens fünf taggleiche OCO-2-Beobachtungen aus neun deutschen MGRS-Tiles enthielten. Die externe Bewertung erfolgte anhand von 2.000 nativen OCO-2-Footprints aus zwei zuvor nicht verwendeten Tiles. Das Ertragsexperiment nutzte Daten von 2019--2022 zum Training und bewertete 48 bayerische NUTS-3-Regionen im Jahr 2024. Das räumlich aggregierte 4-km-U-Net erzielte intern einen RMSE-Wert von 0.0701 und ein $R^2$ von 0.761.
Die externe Bewertung innerhalb des trainierten SIF-Wertebereichs ergab auf nativer Footprint-Ebene einen RMSE-Wert von 0.1434 und ein $R^2$ von 0.349. Diese Leistung war weitgehend mit modernen SIF-Enhancement-Methoden vergleichbar, obwohl Unterschiede in Datensätzen, Untersuchungsgebieten und Validierungsgrundlagen einen direkten Vergleich einschränken.
CNN- und baumbasierte Modelle erzielten auf Footprint-Ebene ähnliche Kennzahlen, sodass Zielkonstruktion und räumliche Bezugsfläche einflussreicher als die Modellarchitektur waren. CNNs konnten jedoch Nachbarschaftsstrukturen und Feldgrenzen besser abbilden. Kulturartenreines SIF erreichte 2024 einen Ertrags-RMSE von 0.9406~$\mathrm{t\,ha^{-1}}$, gegenüber 1.0285--1.0874~$\mathrm{t\,ha^{-1}}$ für gemischte SIF-Footprints. Dies entspricht RMSE-Reduktionen von 8.5\% beziehungsweise 13.5\%. Das Ertragsexperiment bewertet jedoch nur den Einfluss der SIF-Reinheit und nicht die allgemeine Vorhersagegüte. Die Arbeit stellt einen reproduzierbaren Ansatz zur Lückenfüllung und Auflösungsverbesserung von OCO-2-SIF bereit, der keine räumlich vollständige SIF-Ausgangsbeobachtung benötigt. Die erzeugten 20-m-Werte werden als modellbasierte Vorhersagen behandelt und auf nativer OCO-2-Footprint-Ebene validiert. Die Ergebnisse liefern erste Hinweise auf den Nutzen kulturartenspezifischer SIF-Signale für die landwirtschaftliche Ertragsüberwachung und zeigen Schwerpunkte für die Weiterentwicklung des Ansatzes auf.